\section{Zvolená varianta a odůvodnění}
\label{sec:selection-choice}

Na základě analýzy kritérií (sekce~\ref{sec:selection-criteria}), porovnání způsobů nasazení (sekce~\ref{sec:selection-deployment}) a přehledu dostupných modelů (sekce~\ref{sec:selection-models}) tato sekce shrnuje trade-offs jednotlivých řešení a odůvodňuje volbu konkrétní varianty pro implementaci v systému Kelvin.

\subsection{Trade-offs jednotlivých řešení}
\label{subsec:choice-tradeoffs}

Žádné z posuzovaných řešení není ideální ve všech ohledech.
Každá volba s sebou přináší kompromisy mezi kvalitou, cenou, ochranou dat a složitostí nasazení.
Klíčové trade-offs jsou shrnuty níže.

\subsubsection*{Cloudové API s proprietárním modelem}

Varianta využívající cloudové API (například GPT-4o nebo Claude~3.5 Sonnet) nabízí:
\begin{itemize}
    \item \textbf{Výhody:} Nejvyšší kvalita výstupů; nulové hardwarové nároky; jednoduché zahájení provozu; automatické aktualizace modelu; dostatečné kontextové okno.
    \item \textbf{Nevýhody:} Průběžné náklady za token; závislost na externím poskytovateli; potřeba GDPR ošetření (DPA); potenciální výpadky API; výběr modelu omezen na nabídku poskytovatele.
\end{itemize}

\subsubsection*{On-premise nasazení open-source modelu}

Varianta využívající lokálně nasazený open-source model (například Qwen2.5-Coder nebo DeepSeek-Coder-V2):
\begin{itemize}
    \item \textbf{Výhody:} Data zůstávají na školní infrastruktuře (GDPR přímočarý); nulové průběžné náklady za token; úplná kontrola nad modelem i infrastrukturou; možnost fine-tuningu\footnote{Fine-tuning je proces dodatečného doladění předtrénovaného modelu na specifickém datasetu, čímž lze zlepšit jeho výkony pro konkrétní úlohu nebo doménu.} na specifická data.
    \item \textbf{Nevýhody:} Vysoké jednorázové náklady na GPU hardware; nižší kvalita výstupů oproti proprietárním modelům (zejména u menších variant); nutnost správy infrastruktury; zodpovědnost za bezpečnost a aktualizace.
\end{itemize}

\subsubsection*{Hybridní přístup}

Hybridní přístup kombinuje obě varianty -- například cloudové API pro produkční provoz s možností přepnutí na on-premise model v případě potřeby:
\begin{itemize}
    \item \textbf{Výhody:} Flexibilita; ochrana před závislostí na jednom poskytovateli (vendor lock-in).
    \item \textbf{Nevýhody:} Vyšší složitost implementace; nutnost udržovat dvě různé konfigurace.
\end{itemize}

\subsection{Odůvodnění volby konkrétního modelu}
\label{subsec:choice-justification}

Po zvážení všech výše zmíněných faktorů byla pro implementaci v systému Kelvin zvolena varianta využívající \textbf{cloudové API Anthropic s modelem Claude~3.5 Haiku} jako primárním modelem pro produkční provoz, s možností přepnutí na Claude~3.5 Sonnet pro kvalitnější analýzu v případě potřeby.

Tato volba vychází z následujících důvodů:

\begin{enumerate}
    \item \textbf{Kvalita výstupů při analýze kódu.}
    Modely Claude~3.5 od společnosti Anthropic dosahují na úlohách analýzy a code review výsledků srovnatelných s GPT-4o při přitom výhodné cenové politice.
    V manuálním testování na vzorku historických studentských odevzdání z jazyka C++ generovaly modely Claude konzistentně přesné a pedagogicky hodnotné komentáře.

    \item \textbf{Příznivé podmínky pro ochranu dat.}
    Anthropic poskytuje zákazníkům uzavření DPA a výslovně uvádí, že data odesílaná přes API nejsou bez souhlasu zákazníka využívána pro trénování modelů \cite{anthropic_api}.
    To výrazně zjednodušuje GDPR soulad oproti jiným poskytovatelům, kteří tento závazek neformulují tak jasně.
    Anthropic navíc nabízí zpracování dat v EU regionu.

    \item \textbf{Přijatelné náklady.}
    Model Claude~3.5 Haiku je cenově dostupný (0,80 \$ za milion vstupních tokenů) a pro typický objem odevzdání v systému Kelvin (stovky až nízké tisíce odevzdání za semestr) jsou celkové náklady na API v řádu desítek dolarů za semestr -- zanedbatelné ve srovnání s náklady na hardware nutný pro on-premise nasazení.

    \item \textbf{Velké kontextové okno.}
    Kontextové okno 200\,000 tokenů modelů Claude~3.5 je výrazně nadlimitní pro jakékoli odevzdání v systému Kelvin.
    Eliminuje nutnost řešit chunking nebo selektivní kontext i pro rozsáhlejší projekty.

    \item \textbf{Nulové hardwarové nároky.}
    Systém Kelvin nemá k dispozici dedikovaný GPU server vhodný pro on-premise inferenci.
    Cloudové API umožňuje zahájit provoz bez pořizování hardware a je ekonomicky výhodné pro aktuální objem zpracování.

    \item \textbf{Modulární architektura pro budoucí rozšíření.}
    Implementace je navržena tak, aby byl model snadno vyměnitelný.
    Systém komunikuje s modelem přes abstraktní rozhraní, které umožňuje přepnutí na libovolný jiný model kompatibilní s OpenAI API formátem -- ať už jiné cloudové API (OpenAI GPT-4o, Gemini) nebo on-premise model přes \texttt{Ollama}.
    Pokud instituce v budoucnu pořídí GPU server, přechod na lokální open-source model (například Qwen2.5-Coder) bude vyžadovat minimální změny v kódu.
\end{enumerate}

\begin{table}[H]
    \centering
    \resizebox{\textwidth}{!}{%
        \begin{tabular}{lcccc}
            \toprule
            \textbf{Kritérium}              & \textbf{Claude 3.5 Haiku} & \textbf{GPT-4o-mini} & \textbf{Qwen2.5-Coder 7B} & \textbf{Llama 3.1 8B} \\
            \midrule
            Kvalita code review             & $\bullet\bullet\bullet\circ$  & $\bullet\bullet\bullet\circ$  & $\bullet\bullet\circ\circ$   & $\bullet\bullet\circ\circ$ \\
            Ochrana dat (GDPR)              & $\bullet\bullet\bullet\bullet$ & $\bullet\bullet\circ\circ$   & $\bullet\bullet\bullet\bullet$ & $\bullet\bullet\bullet\bullet$ \\
            Náklady (průběžné)              & $\bullet\bullet\bullet\circ$  & $\bullet\bullet\bullet\circ$  & $\bullet\bullet\bullet\bullet$ & $\bullet\bullet\bullet\bullet$ \\
            Náklady (pořizovací)            & $\bullet\bullet\bullet\bullet$ & $\bullet\bullet\bullet\bullet$ & $\circ\circ\circ\circ$       & $\circ\circ\circ\circ$ \\
            Kontextové okno                 & $\bullet\bullet\bullet\bullet$ & $\bullet\bullet\bullet\bullet$ & $\bullet\bullet\bullet\bullet$ & $\bullet\bullet\bullet\bullet$ \\
            Jednoduchost nasazení           & $\bullet\bullet\bullet\bullet$ & $\bullet\bullet\bullet\bullet$ & $\bullet\bullet\circ\circ$   & $\bullet\bullet\circ\circ$ \\
            \midrule
            \textbf{Celkové hodnocení}      & \textbf{Zvoleno}           & Alternativa           & Alternativa (GPU)         & Alternativa (GPU) \\
            \bottomrule
        \end{tabular}
    }
    \caption{Srovnání finálních kandidátů z hlediska klíčových kritérií ($\bullet$ = splňuje, $\circ$ = nesplňuje nebo nevýhoda). Hodnocení vychází z analýzy v sekcích~\ref{sec:selection-criteria}--\ref{sec:selection-models}.}
    \label{tab:final-model-comparison}
\end{table}

\subsubsection*{Zhodnocení zvolené varianty}

Je důležité přiznat limity zvolené varianty.
Cloudové API Anthropic zavádí závislost na externím poskytovateli a průběžné náklady, které mohou při výrazném nárůstu objemu odevzdání být nezanedbatelné.
Z pohledu ochrany dat je nutné udržovat platnou DPA a průběžně sledovat podmínky API (zejména zda a jakým způsobem jsou data zpracovávána).

Z hlediska kvality výstupů je Claude~3.5 Haiku výkonnostně pod úrovní vlajkového modelu Claude~3.5 Sonnet.
Pro případ, kdy je požadována vyšší kvalita analýzy (například pro závěrečné hodnocení rozsáhlého projektu), je systém navržen tak, aby bylo možné jednoduchou konfigurací přepnout na silnější model.

Celkově je zvolená varianta správnou volbou pro aktuální fázi projektu -- nabízí přijatelnou rovnováhu mezi kvalitou, cenou a snadností nasazení bez nutnosti pořizovat specializovaný hardware.
Implementace v systému Kelvin je pak popsána v kapitole~\ref{ch:implementation}.

\endinput
